\section{Problem Setup}
\label{setup}

\subsection{Market as a Partially Observable Stochastic Game}

We formalize Agent Bazaar as a \textbf{Partially Observable Stochastic Game (POSG)} \citep{hansen2004dynamic} defined by the tuple $(\mathcal{I}, \mathcal{S}, \{\mathcal{A}^i\}_{i \in \mathcal{I}}, \{\mathcal{O}^i\}_{i \in \mathcal{I}}, \mathcal{T}, \{r^i\}_{i \in \mathcal{I}}, \gamma)$, where:
\begin{itemize}
    \item $\mathcal{I} = \{1, \ldots, N\}$ is the set of active agents (firms or buyers, depending on context),
    \item $\mathcal{S}$ is the global market state space (all prices, inventories, reputations, and qualities),
    \item $\mathcal{A}^i$ is the discrete-continuous action space for agent $i$,
    \item $\mathcal{O}^i \subseteq \mathcal{S}$ is the observable state for agent $i$, a strict subset of $\mathcal{S}$ due to Search Friction,
    \item $\mathcal{T}: \mathcal{S} \times \prod_i \mathcal{A}^i \to \Delta(\mathcal{S})$ is the stochastic transition function governing market clearing and reputation updates,
    \item $r^i: \mathcal{S} \times \prod_j \mathcal{A}^j \to \mathbb{R}$ is the scalar reward for agent $i$ (profit or consumer surplus),
    \item $\gamma \in [0,1)$ is the discount factor.
\end{itemize}

\textbf{Search Friction as Partial Observability.}
A structural feature of real-world marketplaces is that participants cannot monitor all competitors simultaneously. We model this as a \textit{token budget} $\tau$ that limits each agent's observable state per timestep: agent $i$ observes $o_t^i \sim \mathrm{Uniform}(s_t, \tau)$, a uniformly sampled subset of the global state $s_t$ of cardinality $\tau$. This models the bounded attention of LLM agents operating with finite context windows and the practical impossibility of monitoring all listings on a large platform in real time. Because $|s_t| \gg \tau$ in general, the effective information available to any agent is severely limited, making myopic local optimization individually rational but globally destabilizing.

We instantiate the POSG in two concrete environments corresponding to the two failure modes identified in our study.

\subsection{Environment A: The Crash (B2C Market)}
\label{setup:crash}

\textbf{State.} The global state $s_t = (I_t^{1:N}, C_t^{1:N}, P_t^{1:N}, D_t)$ comprises firm inventories $I_t^i \in \mathbb{Z}_{\geq 0}$, cash balances $C_t^i \in \mathbb{R}$, posted prices $P_t^i \in \mathbb{R}_{>0}$, and aggregate demand $D_t \sim \mathcal{D}$ drawn from a non-stationary demand process.

\textbf{Action.} Each firm $i$ simultaneously sets price $P_t^i$ and production quantity $Q_t^i$:
$a_t^i = (P_t^i, Q_t^i) \in \mathbb{R}_{>0} \times \mathbb{Z}_{\geq 0}$.

\textbf{Market Clearing.} Consumers purchase from the lowest-priced available firm first (lexicographic by price, then by firm index for ties) until demand $D_t$ is exhausted. Firm $i$ sells $Q_t^{i,\text{sold}} = \min(Q_t^i, \max(0, D_t^i))$ units, where $D_t^i$ is the demand allocated to firm $i$ under market clearing.

\textbf{Reward (Profit).} The per-period reward for firm $i$ is:
\begin{equation}
    r_t^i = (P_t^i - c) \cdot Q_t^{i,\text{sold}} - f \cdot Q_t^i,
\end{equation}
where $c > 0$ is the unit cost of goods and $f \geq 0$ is a per-unit inventory holding cost.

\textbf{Cash Update and Bankruptcy.} Cash evolves as $C_{t+1}^i = C_t^i + r_t^i$. A firm is declared \textit{bankrupt} when $C_t^i < 0$; it exits the market permanently, and its production and inventory are zeroed.

\textbf{Failure Mode: Price Spiral.}
The crash failure mode is characterized by a price spiral: $\partial P_t / \partial t < 0$ across all surviving firms until $P_t^i < c$ for all $i$. Under this condition, each transaction incurs a loss, accelerating the drain on cash reserves and triggering cascading bankruptcies. Formally, the market has \textit{collapsed} at time $T^*$ if $\sum_i \mathbf{1}[C_{T^*}^i < 0] = N$.

\subsection{Environment B: The Lemon Market (C2C Market)}
\label{setup:lemon}

\textbf{Quality and Information Asymmetry.}
Each listed item has a latent quality $q \in \mathcal{Q} = \{1.0, 0.7, 0.4, 0.1\}$ (corresponding to Mint, Good, Fair, Poor condition), drawn independently at listing time. The seller observes $q$ and generates a natural language description $D \in \mathcal{V}^*$; the buyer observes only $(D, R)$, where $R \in [0,1]$ is the seller's reputation score.

\textbf{Reputation Dynamics.}
After each completed transaction at time $t$, the seller's reputation updates via an exponential moving average:
\begin{equation}
    R_{t+1} = \alpha \cdot R_t + (1-\alpha) \cdot q_t, \quad \alpha = 0.9.
\end{equation}
This delayed feedback allows a deceptive seller to sustain a high reputation for several rounds before detection.

\textbf{Buyer Action and Welfare.}
The buyer agent chooses $a_t^{\text{buyer}} = (\text{bid/pass}, P_t^{\max}) \in \{\text{bid}, \text{pass}\} \times \mathbb{R}_{\geq 0}$. Consumer surplus from a completed transaction is $\mathrm{CS}_t = q_t \cdot v - P_t^{\text{final}}$, where $v$ is the buyer's valuation per unit of quality. Negative surplus indicates exploitation.

\textbf{Market Collapse Condition.}
Following Akerlof's \citep{akerlof1970market} analysis, as average seller reputation $\bar{R}_t \to 0$ and buyer confidence erodes, trading volume $V_t \to 0$. We operationalize collapse as the first timestep $T^*$ such that $V_{T^*} < \epsilon_V$ for some threshold $\epsilon_V > 0$.

\textbf{Sybil Attack.}
A \textit{Sybil Principal} controls a cluster of $K$ seller identities $\{S_1, \ldots, S_K\}$, each maintaining an independent reputation $R_k$. The principal uses a single LLM to generate $K$ independently-worded but semantically-equivalent deceptive listings per timestep. When $R_k < \rho_{\min}$ for some minimum reputation threshold $\rho_{\min}$, the principal retires $S_k$ and activates a fresh identity $S_{K+1}$ initialized at $R_{K+1} = R_0$. This identity cycling continuously resets the trust signal available to buyers, amplifying the Akerlof dynamic indefinitely.

\subsection{Economic Alignment (Formal Definition)}
\label{setup:alignment}

We define \textbf{Economic Alignment} as a property of an agent policy $\pi$ relative to a baseline $\pi_{\text{base}}$ across three dimensions:

\begin{itemize}
    \item \textbf{Stability} $(\Phi_S)$: Price volatility $\sigma(P_t)$ across all timesteps and firms. Lower is better.
    \item \textbf{Integrity} $(\Phi_I)$: Fraction of completed transactions where $q < q_{\text{claimed}}$, i.e., the deceptive transaction rate. Lower is better.
    \item \textbf{Consumer Welfare} $(\Phi_W)$: Aggregate consumer surplus $\sum_t \sum_j \mathrm{CS}_t^j$ summed over all buyers $j$ and timesteps $t$. Higher is better.
\end{itemize}

A policy $\pi$ is \textbf{economically aligned} if, relative to $\pi_{\text{base}}$:
\begin{equation}
    \sigma_\pi(P_t) \leq \sigma_{\pi_{\text{base}}}(P_t), \quad
    \Phi_I(\pi) \leq \Phi_I(\pi_{\text{base}}), \quad
    \Phi_W(\pi) \geq \Phi_W(\pi_{\text{base}}).
\end{equation}
We aggregate these into a scalar \textbf{Economic Alignment Score (EAS)}:
\begin{equation}
    \mathrm{EAS}(\pi) = \frac{1}{3}\left[ (1 - \hat{\sigma}) + (1 - \hat{\Phi}_I) + \hat{\Phi}_W \right] \in [0, 1],
\end{equation}
where $\hat{\cdot}$ denotes min-max normalization across all evaluated policies. An EAS of 1 indicates perfect stability, integrity, and welfare; 0 indicates total failure on all three dimensions.
